% =====================================================================
% CSE 341: Concepts of Programming Languages
% Project D4 (Part 2): Cumulative AI Usage Journal
% Student: 220104004931 Enes Özkan
% Compile with: tectonic 220104004931_D4_P2.tex
% =====================================================================
\documentclass[11pt,a4paper]{article}

\usepackage[utf8]{inputenc}
\usepackage[T1]{fontenc}
\usepackage{mathpazo}
\usepackage[scaled=0.92]{helvet}
\usepackage[varqu,varl]{zi4}
\usepackage[english]{babel}
\usepackage{amsmath}
\usepackage{microtype}
\usepackage[a4paper,margin=2.25cm,top=2.0cm,bottom=2.15cm]{geometry}

\usepackage[table]{xcolor}
\definecolor{ansblue}{HTML}{1E3A8A}
\definecolor{anscodebg}{HTML}{F9FAFB}
\definecolor{anscaption}{HTML}{6B7280}
\usepackage[colorlinks=true, linkcolor=ansblue, urlcolor=ansblue,
            citecolor=ansblue, pdfborder={0 0 0}]{hyperref}

\usepackage[explicit]{titlesec}
\titleformat{\section}
  {\large\bfseries\sffamily\color{black}}
  {\thesection}{0.65em}{#1}
\titlespacing*{\section}{0pt}{11pt}{5pt}
\setlength{\parskip}{4.5pt plus 1pt minus 1pt}
\setlength{\parindent}{0pt}
\linespread{1.06}
\setcounter{secnumdepth}{0}

\usepackage{array}
\usepackage{tabularx}
\usepackage{ragged2e}
\newcolumntype{Y}{>{\RaggedRight\arraybackslash}X}
\def\UrlBreaks{\do\/\do-\do_\do.\do~}

\newcommand{\code}[1]{\texttt{\detokenize{#1}}}

\usepackage{fancyhdr}
\pagestyle{fancy}
\fancyhf{}
\renewcommand{\headrulewidth}{0pt}
\fancyfoot[L]{\footnotesize\color{anscaption}%
  CSE 341: AI Usage Journal (D4, Part 2 Cumulative) \(\cdot\)
  Student 220104004931 Enes Özkan}
\fancyfoot[R]{\footnotesize\color{anscaption}Page \thepage}

\hypersetup{
  pdftitle  = {AI Usage Journal (D4, Part 2 Cumulative)},
  pdfauthor = {220104004931 Enes Özkan},
  pdfsubject= {CSE 341 Project Part 2},
  pdfkeywords={AI, usage journal, ANS, DSL, type system, interpreter}
}

\newcommand{\field}[2]{\textbf{#1.} #2\par}
\newcommand{\prompt}[1]{%
  \textbf{Prompt.}
  \begin{quote}\small #1\end{quote}
}
\newcommand{\sessionlog}[1]{\textbf{Session log.} #1\par}

\newcommand{\entrygap}{\vspace{8pt}}
\newenvironment{journalentry}{%
  \begin{center}\small\renewcommand{\arraystretch}{1.22}%
  \begin{tabular}{|>{\columncolor{anscodebg}\bfseries}p{3.0cm}|>{\RaggedRight\arraybackslash}p{12.25cm}|}%
  \hline
}{%
  \end{tabular}\end{center}\entrygap
}

\begin{document}

\begin{center}
  {\footnotesize\color{anscaption}CSE 341: Concepts of Programming Languages\par}
  \vspace{3pt}
  {\huge\bfseries\sffamily AI Usage Journal\par}
  \vspace{4pt}
  {\large Deliverable 4 (D4), Part 2 Cumulative Journal\par}
  \vspace{6pt}
  {\footnotesize\color{anscaption}
    Student: 220104004931 Enes Özkan \(\cdot\)
    Spring 2026 \(\cdot\)
    Gebze Technical University\par}
\end{center}
\vspace{7pt}

\field{Submission note}{This Part 2 upload keeps my Part 1 AI Journal entries first, then adds the Part 2 continuation entries. All ChatGPT web-app sessions are from the following project history: \url{https://chatgpt.com/g/g-p-69f85c6705f08191ac4f2f56826649bf-ans/project} (may require ChatGPT login).}

% =====================================================================
\section{Part 1 Journal Entries}

\begin{journalentry}
Entry \# & 1 of 18 \\ \hline
Date & 2026-05-04 \\ \hline
Phase & Design \\ \hline
AI tool & GPT-5.5 Thinking \\ \hline
Goal & Connect the ANS language idea to Sebesta's language evaluation criteria. \\ \hline
Prompt & Which criteria does this DSL prioritize, and which does it knowingly sacrifice according to Sebesta? \\ \hline
Response (key part) & The response argued that ANS should prioritize readability, writability, and reliability because domain terms such as \code{element}, \code{compound}, \code{reaction}, and \code{check = balanced} make chemistry concepts visible in the program. It also said ANS sacrifices generality, orthogonality, and implementation cost. \\ \hline
Accepted & I kept the readability and reliability arguments because they match the project goal: ANS should make chemistry notation clearer and catch reaction mistakes earlier than a general-purpose script. \\ \hline
Rejected / modified & I shortened the wording for D1 and removed broad claims that made ANS sound more complete than the Part 1 lexer and parser actually are. \\ \hline
Errors you caught & The answer was mostly useful, but it treated some future semantic checks as if they already existed. I detected that by comparing it with the D2 implementation, which is still lexer and parser only in Part 1. \\ \hline
Reflection & This was a good example of using AI for framing, not for final truth. It helped me organize Sebesta terminology, but I still had to separate implemented behavior from planned behavior. \\ \hline
\end{journalentry}

\begin{journalentry}
Entry \# & 2 of 18 \\ \hline
Date & 2026-05-04 \\ \hline
Phase & Lexer \\ \hline
AI tool & GPT-5.5 Thinking \\ \hline
Goal & Draft token categories and regular-expression-style patterns for the lexical structure section. \\ \hline
Prompt & Can you explain the list of token categories: identifiers, literals, keywords, operators, separators. Give the regular expression or pattern for each in ANS. \\ \hline
Response (key part) & The response proposed categories for identifiers, number literals, string literals, formula literals, unit literals, keywords, operators, separators, and comments. One formula pattern was similar to \code{(?:[A-Z][a-z]?\d*)+}. \\ \hline
Accepted & I kept the category breakdown and the idea of checking keywords after scanning identifier-shaped text. That matches a normal lexer design and kept the D1 lexical section organized. \\ \hline
Rejected / modified & The first formula pattern was too permissive and did not prove that a formula uses real element symbols. I treated formula validation as a later semantic issue instead of pretending the lexer solves chemistry correctness. \\ \hline
Errors you caught & The AI blurred lexical validity and chemical validity. A token can look like a formula while still not being chemically meaningful, so the report needed to say what the lexer can and cannot know. \\ \hline
Reflection & This was a good example of using AI for framing, not for final truth. It helped me organize Sebesta terminology, but I still had to separate implemented behavior from planned behavior. \\ \hline
\end{journalentry}

\begin{journalentry}
Entry \# & 3 of 18 \\ \hline
Date & 2026-05-05 \\ \hline
Phase & Lexer / Parser \\ \hline
AI tool & GPT-5.5 Thinking \\ \hline
Goal & Clarify how ANS should distinguish chemical formula literals from ordinary identifiers. \\ \hline
Prompt & For ANS, I would treat chemical formulas like H2O, CH4, CO2, and O2 as formula literals. Could you explain how the lexer and parser should distinguish formula literals from identifiers properly? \\ \hline
Response (key part) & The response said the lexer should recognize the shape of the text and the parser should decide the role from grammar position. It gave examples where \code{H2O} appears as a compound formula while lowercase names such as \code{water} remain identifiers. \\ \hline
Accepted & I kept the general principle that parser context matters. It made the design easier to explain because uppercase formula-shaped tokens and lowercase identifiers have different roles in ANS. \\ \hline
Rejected / modified & I did not use the suggested single broad \code{WORD} token. The actual implementation separates \code{FORMULA} and identifier-like tokens because that fit the C lexer and parser better. \\ \hline
Errors you caught & The AI's broad-token suggestion would have moved too much responsibility into the parser for our implementation. I caught this while checking how the recursive-descent parser consumes token kinds. \\ \hline
Reflection & AI gave a useful compiler-design concept, but the implementation still had to be adapted to the actual C code structure. \\ \hline
\end{journalentry}

\begin{journalentry}
Entry \# & 4 of 18 \\ \hline
Date & 2026-05-05 \\ \hline
Phase & Lexer \\ \hline
AI tool & GPT-5.5 Thinking \\ \hline
Goal & Improve the lexical specification for units such as \code{g/mol}, \code{kg/mol}, and \code{kJ/mol}. \\ \hline
Prompt & List the token categories for ANS more completely, including identifiers, literals, keywords, operators, separators, formula literals, and unit literals. Give the regex or pattern for each and think carefully about ambiguity. \\ \hline
Response (key part) & The response suggested a \code{UNIT\_LITERAL} pattern using lookbehind and lookahead assertions, such as a boundary check around \code{kg/mol}, \code{g/mol}, and \code{kJ/mol}. It also recommended longest-match handling before treating \code{/} as division. \\ \hline
Accepted & I kept the longest-match idea for slash units. The current lexer probes \code{g/mol}, \code{kg/mol}, and \code{kJ/mol} before returning the slash operator, which avoids splitting a unit into several tokens. \\ \hline
Rejected / modified & I rejected the lookbehind-based regex because a hand-written C lexer does not use Python-style regex features. The final implementation uses direct string checks instead. \\ \hline
Errors you caught & The AI proposed a technically neat regex that was mismatched with the implementation language. I detected this by checking what would actually be reasonable in plain C. \\ \hline
Reflection & This entry reminded me that a correct-looking formal pattern is not automatically a good implementation plan. The target language matters. \\ \hline
\end{journalentry}

\begin{journalentry}
Entry \# & 5 of 18 \\ \hline
Date & 2026-05-05 \\ \hline
Phase & Design / Parser \\ \hline
AI tool & GPT-5.5 Thinking \\ \hline
Goal & Review the move toward a more C-like ANS syntax and identify parser consequences. \\ \hline
Prompt & We decided to change the syntax to a more C-like type with braces, semicolons, assignments using equals, function declarations, if statements, and line comments. What lexer and parser changes does this imply for ANS? \\ \hline
Response (key part) & The response said C-like syntax would make the grammar more familiar through \code{let}, braces, semicolons, \code{=}, and function-call syntax. It also proposed updated token patterns and grammar fragments. \\ \hline
Accepted & I kept the direction of using semicolons, braces, \code{=}, function calls, and \code{//} comments. Those choices are reflected in the submitted D1 and D2 parser. \\ \hline
Rejected / modified & The AI suggested a \code{compute\_yield} shape that did not match our examples. I did not treat that as final and kept the parser aligned with the D3 programs instead. \\ \hline
Errors you caught & The hallucinated function signature showed that the AI was filling gaps from its own assumptions. I caught it by comparing the suggestion against the example programs we were actually using. \\ \hline
Reflection & Syntax changes ripple through the whole project. This was useful, but only because I checked every suggestion against the examples and implementation. \\ \hline
\end{journalentry}

\begin{journalentry}
Entry \# & 6 of 18 \\ \hline
Date & 2026-05-05 \\ \hline
Phase & Parser \\ \hline
AI tool & GPT-5.5 Thinking \\ \hline
Goal & Generate an unambiguous EBNF grammar and review ambiguity-resolution choices. \\ \hline
Prompt & Can you write the complete EBNF for our DSL ANS. The EBNF must be unambiguous for the subset you implement. Annotate any places where you resolve ambiguity, such as operator precedence and dangling else. \\ \hline
Response (key part) & The response separated expression precedence into additive, multiplicative, unary, and primary rules. It also described reaction terms such as \code{2*O2} and discussed formula-shaped tokens in context. \\ \hline
Accepted & I kept the precedence hierarchy because it directly matches the recursive-descent parser shape. I also kept the idea that stoichiometric coefficients must use \code{*} rather than being written as \code{2O2}. \\ \hline
Rejected / modified & The response incorrectly said there was no dangling-else issue because the subset had no \code{if}/\code{else}. ANS does have \code{if}, \code{elif}, and \code{else}, so I did not use that part. \\ \hline
Errors you caught & The AI contradicted the language features already in our examples. I caught this by checking the grammar against the D3 control-flow programs. \\ \hline
Reflection & This entry is a clear case where AI was useful for structure but unreliable on details. The safest workflow was to compare each grammar claim with a real program. \\ \hline
\end{journalentry}

\begin{journalentry}
Entry \# & 7 of 18 \\ \hline
Date & 2026-05-05 \\ \hline
Phase & Names / Binding / Scope \\ \hline
AI tool & GPT-5.5 Thinking \\ \hline
Goal & Explain ANS identifier rules, binding times, scoping, and lifetime using Sebesta Chapter 5 concepts. \\ \hline
Prompt & Answer for ANS, citing Sebesta Ch. 5: What are the rules for legal identifiers? What bindings happen at compile time vs. run time? Does your language use static scoping or dynamic scoping? Why? What is the lifetime of local variables? \\ \hline
Response (key part) & The response recommended static lexical scoping and stack-dynamic local variables. It separated compile-time bindings such as declaration lookup from run-time bindings such as values produced during execution. \\ \hline
Accepted & I accepted static scoping and stack-dynamic locals because they make functions easier to reason about and match the D1 design. A chemistry calculation should not change because a caller happens to have a local variable with the same name. \\ \hline
Rejected / modified & I removed wording that sounded like a finished interpreter already existed. In Part 1, this is a design commitment and parser-compatible plan, not executed run-time behavior yet. \\ \hline
Errors you caught & The response sometimes blurred D1 design decisions with D2 implementation status. I corrected that distinction in the report. \\ \hline
Reflection & This helped me use Sebesta vocabulary more accurately, especially the difference between binding time, scope, and lifetime. \\ \hline
\end{journalentry}

\begin{journalentry}
Entry \# & 8 of 18 \\ \hline
Date & 2026-05-07 \\ \hline
Phase & Design / Parser \\ \hline
AI tool & GPT-5.5 Thinking \\ \hline
Goal & Ask for a critique of the current ANS design before finishing the implementation. \\ \hline
Prompt & Do you see any issues with the current state of ANS that we will be implementing? \\ \hline
Response (key part) & The response flagged scope creep, contextual keywords, formula literals versus compound identifiers, and an unclear \code{compute\_yield} signature. It recommended that reaction species allow both formula literals and compound identifiers. \\ \hline
Accepted & I accepted the warning about scope creep and the need to handle both formula literals and identifiers as reaction species. Those are real parser risks in a chemistry DSL. \\ \hline
Rejected / modified & I did not accept every proposed semantic feature, because Part 1 only needed the lexer and parser. Some suggestions were moved into the Part 2 plan instead. \\ \hline
Errors you caught & The AI treated several type-checker and interpreter concerns as urgent Part 1 work. I caught that by checking the handout, which only requires lexer and parser for D2 Part 1. \\ \hline
Reflection & A critique prompt was more useful than asking for code. It helped identify design risks without pretending that all of them had to be solved immediately. \\ \hline
\end{journalentry}

\begin{journalentry}
Entry \# & 9 of 18 \\ \hline
Date & 2026-05-07 \\ \hline
Phase & Lexer \\ \hline
AI tool & GPT-5.5 Thinking \\ \hline
Goal & Generate a lexer from the ANS EBNF and design specification. \\ \hline
Prompt & I want you to build a working lexer according to our EBNF for our programming language we're creating named ANS. Make sure to follow the file I provided to you accordingly and properly. \\ \hline
Response (key part) & The response created a Python lexer with token classes, regex scanning, keyword handling, and unit handling. It kept chemistry words such as \code{element}, \code{compound}, and \code{reaction} contextual rather than hard keywords. \\ \hline
Accepted & I kept the keyword-classification idea, especially the split between hard programming keywords and contextual chemistry words. \\ \hline
Rejected / modified & I rejected the actual generated code because it was Python and our implementation is in C. The prompt was underspecified because I did not state the implementation language. \\ \hline
Errors you caught & Wrong target language. If accepted blindly, it would not fit the project structure or Makefile. \\ \hline
Reflection & This taught me to include hard constraints in implementation prompts. AI often chooses Python for scanners unless the target language is explicit. \\ \hline
\end{journalentry}

\begin{journalentry}
Entry \# & 10 of 18 \\ \hline
Date & 2026-05-07 \\ \hline
Phase & Lexer \\ \hline
AI tool & GPT-5.5 Thinking \\ \hline
Goal & Translate the lexer approach into C and compare it with the project structure. \\ \hline
Prompt & The code needs to be in C, and it should follow the ANS lexer and parser design we discussed: C-like syntax, formula literals, string and char literals, slash units, hard keywords, and line comments. \\ \hline
Response (key part) & The response rebuilt the lexer idea in C with token enums, scanner helper functions, keyword filtering, and direct checks for compound units such as \code{kg/mol}, \code{kJ/mol}, and \code{g/mol}. \\ \hline
Accepted & I accepted the manual-scanning approach and the direct unit checks because they match a hand-written C lexer. \\ \hline
Rejected / modified & I adapted the generated structure to the actual repository files and error format. I did not keep any claim that comment handling had extra build-time options, since the current lexer simply skips \code{//} comments. \\ \hline
Errors you caught & The AI's code shape did not automatically match our existing file layout and naming. I caught that during integration with \code{D2/include} and \code{D2/src}. \\ \hline
Reflection & AI was useful for translating the idea into C, but repository integration still required manual checking against headers, token names, and tests. \\ \hline
\end{journalentry}

\begin{journalentry}
Entry \# & 11 of 18 [EXPERIMENT E1] \\ \hline
Date & 2026-05-08 \\ \hline
Phase & Parser \\ \hline
AI tool & GPT-5.5 Thinking \\ \hline
Goal & \textbf{Experiment E1:} ask AI to generate an EBNF grammar from a one-paragraph language description. \\ \hline
Prompt & I am designing a domain-specific programming language called ANS (Atomic Notation Script) for chemistry. It has C-like syntax with statements ending in semicolons and blocks in curly braces. The language supports standard constructs like 'let' variable declarations, 'if/elif/else' conditionals, 'for ... in' loops, and 'function' definitions. It also has domain-specific constructs: users can declare an 'element' with properties like atomic\_mass and valence, a 'compound' with a chemical formula literal like H2O or CH4, and a 'reaction' that defines reactants, products, a 'check = balanced' constraint, and energy changes such as 'exothermic 890.3 kJ/mol'. The expression language supports standard arithmetic (+, -, *, /) and comparison operators. Please generate a complete EBNF grammar for this language. \\ \hline
Response (key part) & The response produced a full grammar with statements, declarations, formulas, reactions, and expression precedence. It also added primitive type names such as \code{number}, \code{string}, and \code{bool}, plus extra domain details not requested in the prompt. \\ \hline
Accepted & I accepted the general hierarchy and the precedence structure, because the expression grammar was close to what a recursive-descent parser needs. \\ \hline
Rejected / modified & I rejected the invented primitive type declarations, nested formula forms, and hardcoded chemical states. Our Part 1 language has literals, typed domain parameters, and flat formula literals, but not all of those extra grammar features. \\ \hline
Errors you caught & The AI overfilled gaps in the one-paragraph description with generic language-design assumptions. I detected this by comparing the generated grammar to D1 and the D3 examples. \\ \hline
Reflection & E1 showed that a short prompt is not enough for a precise grammar. AI handles the common structure well, but it invents details unless constraints are made explicit. \\ \hline
\end{journalentry}

\begin{journalentry}
Entry \# & 12 of 18 \\ \hline
Date & 2026-05-08 \\ \hline
Phase & Parser \\ \hline
AI tool & GPT-5.5 Thinking \\ \hline
Goal & Refine the E1 grammar after identifying missing constraints and ambiguous syntax. \\ \hline
Prompt & Good start, but fix these constraints: in reactions require a * between a coefficient and a species, as in 2 * O2; do not invent primitive type declarations; function parameters may use domain types such as element, compound, and reaction; chemical formulas are flat; all top-level declarations end with semicolons; units like g/mol, kg/mol, and kJ/mol are unit tokens; do not invent syntax that was not stated. \\ \hline
Response (key part) & The response revised the grammar so reaction terms use an optional coefficient followed by \code{*} and a species reference, top-level declarations end with semicolons, formulas stay flat, and unit syntax is represented separately from division. \\ \hline
Accepted & I accepted the corrections that matched the implementation direction: explicit \code{*} in reaction terms, flat formulas, semicolon-terminated top-level declarations, and separate treatment for slash units. \\ \hline
Rejected / modified & I still reviewed the final grammar manually before using it, because the parser had to match exact C implementation functions and the actual D3 programs. \\ \hline
Errors you caught & The first E1 output had ambiguity and invented features. The corrected prompt reduced those problems, which showed that negative constraints are necessary when asking AI for formal syntax. \\ \hline
Reflection & The useful part was not just the final grammar. The iteration exposed exactly which language details had to be decided before the parser could be written safely. \\ \hline
\end{journalentry}

% =====================================================================
\section{Part 2 Continuation Entries}

\begin{journalentry}
Entry \# & 13 of 18 \\ \hline
Date & 2026-05-19 \\ \hline
Phase & Design \\ \hline
AI tool & GPT-5.5 Thinking \\ \hline
Goal & Write operational semantics for two non-trivial ANS chemistry DSL constructs: formula iteration and reaction balance checking. \\ \hline
Prompt & We're doing Part 2 of our CSE341 language project. For D1 section 4.4, we need to write operational semantics for two non-trivial ANS constructs. Our language is a chemistry DSL. Can you help us write the semantics section for: 1. looping over a formula, like for (let element in formula) 2. checking whether a reaction is balanced with check = balanced \\ \hline
Response (key part) & The response drafted a D1-ready operational semantics section using environments \(\Gamma\), \(\rho\), and \(\sigma\). It defined formula values such as \(H_2O \Downarrow \{ H \mapsto 2, O \mapsto 1 \}\), gave a formal loop rule for \(\langle \text{for (let element in e) B}, \rho, \sigma_0, \Gamma \rangle \Downarrow \sigma_n\), and defined \code{check = balanced} by comparing \(atoms(R) = atoms(P)\). It also stated that \code{check = balanced} ``does not try to repair or auto-balance the reaction.'' \\ \hline
Accepted & I accepted the environment definitions (\(\Gamma\), \(\rho\), \(\sigma\)) and the formal representation of formulas as a map from element symbols to counts (e.g. \(H_2O \Downarrow \{H \mapsto 2, O \mapsto 1\}\)), as well as the atom conservation constraint for checking if reactions are balanced. \\ \hline
Rejected / modified & I rejected the logical inference-rule format for D1 because our course requires Sebesta-style operational semantics, which are based on abstract-machine state transitions. We also modified the looping definition to evaluate general expressions to formulas rather than assuming only literal formula variables. \\ \hline
Errors you caught & The AI's initial inference rules assumed a static set of elements in the formula and didn't define iteration order or local variable scope cleanup, which would cause name shadowing bugs in our interpreter. \\ \hline
Reflection & I learned that formal logic-style operational semantics are clean but hard to map directly to an interpreter implementation in C. Switching to Sebesta-style operational semantics makes the translation to C code much more natural. \\ \hline
\end{journalentry}

\begin{journalentry}
Entry \# & 14 of 18 \\ \hline
Date & 2026-05-19 \\ \hline
Phase & Design \\ \hline
AI tool & GPT-5.5 Thinking \\ \hline
Goal & Revise the ANS operational semantics section so it follows Sebesta Section 3.5 style more closely. \\ \hline
Prompt & The Semantics needs to be like Sebesta Style as you can see through the books 3.5 section \\ \hline
Response (key part) & The response replaced formal inference rules with Sebesta-style operational descriptions using abstract-machine operations such as \code{evaluate(E)}, \code{formula\_elements(F)}, \code{reaction\_atoms(R)}, and \code{semantic\_error(M)}. It described \code{for (let element in formula)} as lower-level loop operations and \code{check = balanced} as comparing \code{left\_atoms} and \code{right\_atoms} before calling \code{bind\_reaction(N, reaction\_value)}. \\ \hline
Accepted & We accepted the use of abstract machine state modifications and helper functions (like \code{evaluate}, \code{formula\_elements}, \code{reaction\_atoms}, and \code{semantic\_error}) to specify semantics, as this maps directly to how our compiler front-end operates. \\ \hline
Rejected / modified & The AI's list of helper functions was initially unstructured. We restructured the description to explicitly define each helper function before using it in the operational rules, ensuring the machine state remains formal. \\ \hline
Errors you caught & The AI used \code{formula\_elements(F)} like a mathematical set without specifying a deterministic iteration order. This is a semantic gap for interpreter execution. We caught this and defined the iteration order to be the left-to-right order of appearance in the formula. \\ \hline
Reflection & Sebesta's style of operational semantics provides a clear blueprint for interpreter implementation because it describes semantics using low-level, step-by-step algorithms on an abstract machine. \\ \hline
\end{journalentry}

\begin{journalentry}
Entry \# & 15 of 18 \\ \hline
Date & 2026-05-22 \\ \hline
Phase & Type Checker \\ \hline
AI tool & GPT-5.5 Thinking \\ \hline
Goal & Decide how to implement ANS static lexical scoping and symbol-table lookup in C11. \\ \hline
Prompt & Okay, next up is the symbol table for the type checker. We need static lexical scoping (like Sebesta Ch 5) so local loop variables don't overwrite global element declarations. In C, is it fine to just use a linked list of Env structs where each one points to its parent scope, or is that going to be way too slow for lookups? \\ \hline
Response (key part) & The response recommended using \code{Env} structs with a \code{parent} pointer and local linked-list symbols, explaining that insertion should use only the current scope while lookup walks outward. Key code included \code{Symbol *env\_lookup(Env *env, const char *name)} and \code{bool env\_define(Env *env, const char *name, SymbolKind kind, Type type)}. It also noted that this is fast enough for ANS Part 2 and can later be upgraded to hash buckets per scope. \\ \hline
Accepted & We accepted the lexical environment structure using a parent pointer linked list. It is very simple to implement in C, matches Sebesta's static scoping rules, and keeps local variables isolated in nested scopes. \\ \hline
Rejected / modified & The AI suggested that we allow variable shadowing for all symbols including elements and compounds. We rejected this and modified the implementation to forbid shadowing global compound and element names, as shadowing these in a chemistry DSL is highly confusing. \\ \hline
Errors you caught & The AI suggested that lookup complexity was a non-issue because ``compilers rarely have deep nesting.'' While true for standard scripts, it did not account for the fact that a naive linear search on large codebases could scale poorly, though it is fine for our course project's constraints. \\ \hline
Reflection & Designing scoping environments is simple if we treat scope as a first-class linked structure. Restricting shadowing is a design decision that improves reliability for domain-specific languages. \\ \hline
\end{journalentry}

\begin{journalentry}
Entry \# & 16 of 18 \\ \hline
Date & 2026-05-22 \\ \hline
Phase & Interpreter \\ \hline
AI tool & GPT-5.5 Thinking \\ \hline
Goal & Fix memory ownership problems in the C11 ANS interpreter caused by shallow-copying pointers inside tagged-union values. \\ \hline
Prompt & We got the interpreter mostly working, but I'm getting destroyed by Valgrind right now. When I assign a string or a formula to a variable, my tagged union just copies the pointer. So when the local scope ends, I either get double-free errors or memory leaks. How do I handle deep copying in a recursive AST walker without leaking memory everywhere? \\ \hline
Response (key part) & The response recommended making ownership explicit: ``Every Value returned by eval\_expr owns its heap data'' and ``Every variable binding owns its stored Value.'' It introduced \code{value\_free}, \code{value\_clone}, \code{value\_move}, and \code{value\_replace} so strings and formulas are deep-copied or moved safely. It also recommended initializing temporaries to \code{VAL\_VOID} and using a single \code{cleanup:} block in recursive evaluator functions. \\ \hline
Accepted & We accepted the rule of explicit value ownership (every returned \code{Value} owns its heap memory) and the C patterns of \code{value\_clone}, \code{value\_move}, and \code{value\_replace}, along with using a single \code{cleanup:} label in recursive evaluator functions. \\ \hline
Rejected / modified & We rejected the suggestion of using C++ style smart pointers or complex reference counting, keeping it strictly to standard C allocations to avoid complicating the interpreter structure. \\ \hline
Errors you caught & The AI's initial implementation of value cloning did not deeply clone formula parts, which would have led to shared formula arrays and double-frees during variable re-assignments. I detected this by checking the formula clone logic and added nested allocations. \\ \hline
Reflection & Managing memory in a C interpreter requires a clear ownership model. Initializing all temporaries to a void state and freeing them in a unified cleanup block is the cleanest way to prevent resource leaks. \\ \hline
\end{journalentry}

\begin{journalentry}
Entry \# & 17 of 18 \\ \hline
Date & 2026-05-22 \\ \hline
Phase & Type Checker \\ \hline
AI tool & GPT-5.5 Thinking \\ \hline
Goal & Design the recursive C11 type-checker pass for ANS expressions with no implicit coercion. \\ \hline
Prompt & I need to write the actual type checker pass that runs before the interpreter. Since ANS has zero implicit coercion, we have to catch stuff like 5 mol + 2 g. How do I write a recursive C function that walks an expression AST and verifies types without actually executing the code? Do I just use a giant switch statement that returns a Type struct and throws an error if binary operator types don't match? \\ \hline
Response (key part) & The response recommended implementing \code{bool check\_expr(TypeChecker *tc, const Expr *expr, Type *out, Error *err)}, where the type checker returns \code{Type} instead of runtime \code{Value}. It used a recursive switch over expression kinds and showed strict type equality with units, including rejecting \code{5 mol + 2 g} because \code{Quantity<mol>} and \code{Quantity<g>} are not equivalent. It also separated \code{check\_lvalue} from \code{check\_expr} so function calls cannot be assignment targets. \\ \hline
Accepted & We accepted the recursive AST-walking \code{check\_expr} design using a type representation that includes the base type and unit, as well as separating \code{check\_lvalue} from \code{check\_expr}. \\ \hline
Rejected / modified & We rejected the generic representation of lists as arbitrary objects, modifying the type system to enforce strict list homogeneity where lists must contain a single element type. \\ \hline
Errors you caught & The AI's code for member access didn't verify if the accessed field existed on the specific type before returning its type (e.g. it allowed any field name on element objects). I caught this and hardcoded validation rules for element and compound fields. \\ \hline
Reflection & Static type checking in C can be clean if we represent types as static structures and perform strict structure checks recursively. Separating type checking from execution is key to static safety. \\ \hline
\end{journalentry}

\begin{journalentry}
Entry \# & 18 of 18 \\ \hline
Date & 2026-05-22 \\ \hline
Phase & Diagnostics \\ \hline
AI tool & GPT-5.5 Thinking \\ \hline
Goal & Design compiler diagnostics infrastructure to unwind recursive AST checker stack cleanly without calling exit(1). \\ \hline
Prompt & our grading script uses make test which feeds bad .ans files to our program and expects it to exit with code 1 and a clean <file>:<line>: error: message. Right now, if my type checker finds an error deep in a recursive AST function, I just call exit(1). I know that's terrible C practice, but is it fine for a toy compiler? Or should I be setting a global g\_had\_error flag and letting the whole call stack unwind naturally? \\ \hline
Response (key part) & The response advised against calling \code{exit(1)} deep inside the type checker because it causes memory leaks of AST nodes. It recommended using a \code{Diagnostics} struct to record errors, returning \code{false} from recursive checker functions to unwind the stack, and freeing all AST memory at the top level before exiting. \\ \hline
Accepted & We accepted the pattern of returning \code{bool} in recursive functions and passing a diagnostics reporter struct to handle error collection while allowing the stack to unwind and clean up memory. \\ \hline
Rejected / modified & We rejected using a global variable for tracking error flags. We instead encapsulated the state in a \code{Diagnostics} context passed down to ensure thread-safety and reusability of the compiler library. \\ \hline
Errors you caught & The AI's code did not handle the case where a child node check fails but we still need to avoid evaluating subsequent children (e.g. evaluating the right side of a binary operator when the left side failed). I added early returns on child check failures. \\ \hline
Reflection & Graceful error unwinding in recursive compilers is essential for preventing leaks. It is always better to design cleanup paths early rather than relying on the OS to clean up via abrupt exit calls. \\ \hline
\end{journalentry}

\end{document}
